% Options for packages loaded elsewhere
\PassOptionsToPackage{unicode}{hyperref}
\PassOptionsToPackage{hyphens}{url}
\documentclass[
  10pt,
]{article}
\usepackage{xcolor}
\usepackage[margin=25mm]{geometry}
\usepackage{amsmath,amssymb}
\setcounter{secnumdepth}{-\maxdimen} % remove section numbering
\usepackage{iftex}
\ifPDFTeX
  \usepackage[T1]{fontenc}
  \usepackage[utf8]{inputenc}
  \usepackage{textcomp} % provide euro and other symbols
\else % if luatex or xetex
  \usepackage{unicode-math} % this also loads fontspec
  \defaultfontfeatures{Scale=MatchLowercase}
  \defaultfontfeatures[\rmfamily]{Ligatures=TeX,Scale=1}
\fi
\usepackage{lmodern}
\ifPDFTeX\else
  % xetex/luatex font selection
\fi
% Use upquote if available, for straight quotes in verbatim environments
\IfFileExists{upquote.sty}{\usepackage{upquote}}{}
\IfFileExists{microtype.sty}{% use microtype if available
  \usepackage[]{microtype}
  \UseMicrotypeSet[protrusion]{basicmath} % disable protrusion for tt fonts
}{}
\makeatletter
\@ifundefined{KOMAClassName}{% if non-KOMA class
  \IfFileExists{parskip.sty}{%
    \usepackage{parskip}
  }{% else
    \setlength{\parindent}{0pt}
    \setlength{\parskip}{6pt plus 2pt minus 1pt}}
}{% if KOMA class
  \KOMAoptions{parskip=half}}
\makeatother
\setlength{\emergencystretch}{3em} % prevent overfull lines
\providecommand{\tightlist}{%
  \setlength{\itemsep}{0pt}\setlength{\parskip}{0pt}}
\usepackage{bookmark}
\IfFileExists{xurl.sty}{\usepackage{xurl}}{} % add URL line breaks if available
\urlstyle{same}
\hypersetup{
  pdftitle={SP-10: a bipartite degree-four theorem and an exact duality obstruction},
  pdfauthor={Sidney Holden},
  hidelinks,
  pdfcreator={LaTeX via pandoc}}

\title{SP-10: a bipartite degree-four theorem and an exact duality
obstruction}
\author{Sidney Holden}
\date{14 September 2026}

\usepackage{fvextra,xurl}
\DefineVerbatimEnvironment{verbatim}{Verbatim}{breaklines=true,fontsize=\small}
\begin{document}
\maketitle

\textbf{Affiliation:} Center for Computational Biology, Flatiron
Institute, Simons Foundation.

\textbf{Submission disposition:} Partially resolved. Sections 2-5 pass
for the polynomial complement representation, partial-2-tree
construction, bipartite degree-four theorem, and support obstruction.
The universal graph complement conjecture remains open. Section 6
inherited claims are outside this fresh audit.

\href{https://github.com/ajt60gaibb/OpenProblemsInNLA/blob/main/reviews/2026-09-14-further-submissions/SP-10-review.md}{Independent
informal AI-agent review}. ChatGPT assistance is disclosed. No Lean
verification, external human peer review or formal verification is
claimed. This dated notice supersedes historical statements about
pending review; inherited claims are accepted only within the linked
review scope.

\textbf{Status: partial results. The universal graph complement
conjecture is neither proved nor disproved in this report.}

The required inequality is \[
 \operatorname{mr}(G)+\operatorname{mr}(\overline G)\le n+2.
\] Here the diagonal entries of the admissible real symmetric matrices
are free, while every edge must receive a nonzero off-diagonal entry and
every nonedge must receive zero. A rank estimate without these support
conditions does not answer the problem.

This delivery preserves the available preceding archives, and, when
present in the runtime, the files from the interrupted fourth
continuation. The present report supplies complete proofs of the
specific additional statements below. It does not describe them as a
complete solution or claim priority over the literature. The stored
second-round report contains a withdrawn degeneracy premise; the
third-round correction remains in force.

\subsection{1. Definitions and
dependencies}\label{definitions-and-dependencies}

Let \(\mathcal S(G)\) be the set of real symmetric matrices with the
prescribed off-diagonal support. Let \(\operatorname{mr}_+(G)\) denote
its minimum positive-semidefinite rank. A faithful Gram representation
of \(G\) consists of vectors \(x_v\) such that, for distinct vertices,
\[
 \langle x_u,x_v\rangle\ne0\quad\Longleftrightarrow\quad uv\in E(G).
\] Write \(\rho(G)\) for the minimum dimension when every vector must be
nonzero. Isolated vertices may use zero vectors for ordinary minimum PSD
rank, but not for \(\rho\). Consequently
\(\rho(G)=\operatorname{mr}_+(G)+i(G)\), where \(i(G)\) counts the
isolated vertices. In particular \[
 \operatorname{mr}(G)\le\operatorname{mr}_+(G)\le\rho(G).
\]

The polynomial and duality theorems in this report are elementary and
independent of Hall's theorem. The bipartite degree-four conclusion uses
the standard minor-monotonicity of the positive-semidefinite
Strong-Arnold parameter \(\nu\), not Hall's new minimum-degree result.
Specifically, the only parameter facts needed are \[
 H\text{ a minor of }G\ \Longrightarrow\ \nu(H)\le\nu(G),\qquad
 \nu(G)\le n-\operatorname{mr}_+(G),\qquad \nu(K_t)=t-1\quad(t\ge2).
\] The last equality follows directly from the all-ones matrix and
positive rank. These standard definitions and minor monotonicity are
also recorded in Hall's primary manuscript {[}1, Section 3{]}. The
dependency is explicit: no independent proof of general minor
monotonicity is claimed here.

The elementary graph arguments below use Menger's theorem in its
familiar three-fan form: in a 3-connected graph a vertex has three
internally disjoint paths to a cycle not containing it, with distinct
ends on the cycle. All remaining steps of the partial-2-tree reduction
are included.

\subsection{2. An explicit polynomial complement
representation}\label{an-explicit-polynomial-complement-representation}

\textbf{Theorem 2.1.} Let \(G\) be bipartite with specified parts
\(A=\{a_1,\ldots,a_p\}\) and \(B\). Suppose every vertex of \(B\) has
degree at most \(d\). If the graph has at least one vertex, then \[
 \boxed{\rho(\overline G)\le d+1.}
\] When one part is empty, the graph is edgeless and dimension one
suffices directly.

\textbf{Proof.} Choose distinct positive integers \(t_1,\ldots,t_p\),
and assign to \(a_i\) the vector \[
 u_i=(1,t_i,t_i^2,\ldots,t_i^d)\in\mathbb Z^{d+1}.
\] For \(b\in B\), let \(k_b=|N_G(b)|\) and define the monic polynomial
\[
 p_b(t)=t^{d-k_b}\prod_{a_i\in N_G(b)}(t-t_i).
\] Let \(v_b\in\mathbb Z^{d+1}\) be its coefficient vector, in
increasing powers of \(t\). Then \[
 \langle u_i,v_b\rangle=p_b(t_i).
\] Since every \(t_i\) is positive and distinct, this is zero exactly
when \(a_i\in N_G(b)\). These are exactly the required cross-part zeros
for \(\overline G\).

For two different vertices of \(A\), \[
 \langle u_i,u_j\rangle=\sum_{h=0}^d(t_i t_j)^h>0.
\] For a coefficient in degree \(h\), the sign of its nonzero value in
every \(p_b\) is \((-1)^{d-h}\): the nonzero roots are positive, and the
remaining roots are zero. Thus the products of corresponding
coefficients of any two \(p_b\)'s are nonnegative. Their leading
coefficients are both one. Therefore \[
 \langle v_b,v_c\rangle\ge1
\] for distinct \(b,c\). All within-part edges of the complement are
present. Every vector is nonzero, since \(u_i\) has first coordinate one
and \(v_b\) has last coordinate one. This proves the theorem. The
argument includes isolated vertices and degrees strictly less than
\(d\). \(\square\)

The factor \(t^{d-k_b}\) is important. Without it, polynomials of
different degrees would not have the same coefficient-sign pattern, and
the within-\(B\) nonzero products would not be justified by this
argument.

The construction is deterministic, integral, and involves no
small-singular-value tolerance. It supplies a complement witness, not a
claim that the displayed dimension is always minimal.

\subsection{3. A four-dimensional construction for partial
2-trees}\label{a-four-dimensional-construction-for-partial-2-trees}

\textbf{Theorem 3.1.} Suppose \(G\) is a subgraph of a graph \(F\) with
an ordering in which every vertex's earlier neighbors in \(F\) form a
clique of size at most two. Then \[
 \boxed{\rho(\overline G)\le4.}
\]

\textbf{Proof.} In the given ordering, construct vectors in
\(\mathbb R^4\) with orthogonality exactly at the edges of \(G\).
Maintain the following invariants among the assigned vectors: every
vector is nonzero; every two distinct vectors are independent; and
whenever \(ij\) is an edge of \(F\), the vectors \(x_i,x_j,x_z\) are
independent for every other assigned vertex \(z\).

At the next vertex \(v\), let \(S\) be its set of earlier neighbors in
\(F\), and \(T\subseteq S\) its earlier neighbors in \(G\). The
permitted space is \[
 L=\operatorname{span}\{x_t:t\in T\}^{\perp}.
\] It has dimension \(4-|T|\ge2\). An earlier nonneighbor \(z\notin T\)
must have nonzero product with the new vector. The forbidden hyperplane
\(x_z^\perp\) cannot contain \(L\). Otherwise \(x_z\) lies in the span
of the \(T\)-vectors. For \(|T|\le1\) this violates pairwise
independence; for \(|T|=2\), the pair \(T=S\) is an edge of \(F\), and
the triple-independence invariant gives the contradiction.

To maintain pairwise independence, exclude the line of each earlier
vector. To maintain the invariant for an old edge \(ij\in E(F)\),
exclude \(\operatorname{span}(x_i,x_j)\). If \(|T|\le1\), the permitted
space has dimension at least three and cannot be contained in this
plane. If \(|T|=2\), containment would be equality of two planes. If the
pairs \(\{i,j\}\) and \(T\) intersect, equality is impossible because a
nonzero vector cannot belong both to a space and its orthogonal
complement. If they are disjoint, equality implies that both \(x_i,x_j\)
are perpendicular to both vectors of \(T\). The established support
conditions then give all four cross edges in \(G\), hence in \(F\).
Together with the edge \(ij\) and the edge on \(T\), this would make a
\(K_4\) subgraph of \(F\). Such a subgraph is impossible: its latest
vertex would have three earlier neighbors.

There are also new edges \(vw\) of \(F\), where \(w\in S\). To ensure
independence of \(x_v,x_w,x_z\), exclude
\(\operatorname{span}(x_w,x_z)\) for every earlier \(z\ne w\). Again a
permitted space of dimension at least three cannot be contained in such
a plane. In the remaining case \(|T|=2\), we have \(T=S\), so
\(L\subseteq x_w^\perp\); this plane cannot equal a plane containing the
nonzero vector \(x_w\).

Every excluded subspace is therefore proper in \(L\). A finite union of
proper linear subspaces cannot cover \(L\). Choose a vector outside
their union. This satisfies all support requirements and all induction
invariants. Starting from the empty assignment completes the
construction. Since all constraints and previous vectors can be
rational, rational choices are possible at each stage. \(\square\)

For completeness, the graph condition used here covers every graph with
no \(K_4\) minor. One can obtain the required ordering by repeatedly
deleting a vertex of degree at most two and filling its neighbor pair
with an edge. When the degree is two, the filled graph after deletion is
a minor of the preceding graph: contract one incident edge. Hence it
still has no \(K_4\) minor.

Here is a proof that a \(K_4\)-minor-free graph always has an available
vertex of degree at most two. A 3-connected graph has a \(K_4\) minor:
choose a vertex and a cycle avoiding it, use three disjoint paths from
the vertex to the cycle, and contract the cycle arcs and the paths. In a
2-connected \(K_4\)-minor-free graph of order at least four, there is
consequently a separating pair \(x,y\). For each component after
deleting that pair, take its torso with \(x,y\) and add the edge \(xy\).
Each torso is a smaller 2-connected minor, using a path through another
component to supply the added edge. Induction shows that each torso is
either a triangle or has two nonadjacent degree-two vertices. Since
\(x,y\) are adjacent in a torso, an interior degree-two vertex is
available in either case. Taking such vertices from two components gives
two nonadjacent degree-two vertices in the original graph. Triangles are
the base case. Finally, in a graph with cut vertices, an end block has
an interior vertex of degree at most two by this assertion, or is a
single edge. Disconnected graphs are handled componentwise.

Reversing the deletion-and-fill order gives exactly the
clique-of-size-at-most-two ordering used in Theorem 3.1. Thus the
four-dimensional representation follows without an unproved
generic-position assertion.

\subsection{4. The bipartite degree-four
conclusion}\label{the-bipartite-degree-four-conclusion}

\textbf{Theorem 4.1.} Let \(G\) be a bipartite graph for which one
specified part has all degrees at most four. Then \[
 \boxed{\operatorname{mr}_+(G)+\operatorname{mr}_+(\overline G)\le n+2.}
\] In particular, \(G\) satisfies SP-10. There is no bound on its order
or on the degrees in the other part.

\textbf{Proof.} If \(G\) contains a \(K_4\) minor, minor monotonicity of
\(\nu\) gives \(n-\operatorname{mr}_+(G)\ge3\). Theorem 2.1 gives
\(\operatorname{mr}_+(\overline G)\le5\). Their sum is at most
\(n-3+5=n+2\).

If \(G\) has no \(K_4\) minor but contains a cycle, it has a \(K_3\)
minor, giving \(\operatorname{mr}_+(G)\le n-2\). Theorem 3.1 and the
accompanying graph reduction give
\(\operatorname{mr}_+(\overline G)\le4\). Again the sum is at most
\(n+2\).

If \(G\) is a nonempty forest with \(c\) components, its component
Laplacians give a PSD matrix of rank \(n-c\le n-1\). The complement of
any forest has a nonzero Gram representation in dimension at most three.
To see this directly, order each tree from its root so that each vertex
has at most one earlier neighbor. In \(\mathbb R^3\), choose the new
vector perpendicular to that neighbor, when it exists, and avoid zero
products with all earlier nonneighbors and parallelism with every
earlier vector. The permitted plane has dimension two; pairwise
nonparallelism makes every unwanted orthogonality a proper line, and the
parallelism conditions also exclude at most lines. A finite union of
lines does not cover the plane. Roots impose no perpendicularity and can
be chosen in the whole space. Thus the rank sum is at most
\(n-1+3=n+2\). For the empty vertex set both ranks are zero. \(\square\)

The same proof covers one-side degree at most three with only the
cycle/forest distinction. At degree four, the
\(K_4\)-minor/partial-2-tree distinction is the additional ingredient.
Degree five is not covered by simply reusing this argument: the
polynomial bound becomes six, and a \(K_4\) minor saves only three
ranks.

\subsection{5. Why the rank-two SDP shortcut is not a
proof}\label{why-the-rank-two-sdp-shortcut-is-not-a-proof}

A natural proposed route is to optimize a rank-two objective over a
sparse correlation-matrix spectrahedron. A primal matrix \(X\) and dual
slack \(S\) satisfy \(XS=0\), so their ranks sum to at most \(n\).
Adding a rank-two objective matrix \(C\) to \(S\) gives \(B=S+C\), with
rank at most \(\operatorname{rank}S+2\). The dual constraints can make
\(B\) supported on the complement. This yields the attractive rank count
\[
 \operatorname{rank}X+\operatorname{rank}B\le n+2.
\] But it proves GCC only if \textbf{both} matrices have the exact
required support. Genericity cannot simply be asserted to repair this
problem.

\textbf{Theorem 5.1 (open-set support failure).} For the path
\(1-2-3-4\), consider \[
 \max\{\langle C,X\rangle:X\succeq0,\ X_{ii}=1,\ X_{13}=X_{14}=X_{24}=0\}.
\] Write \(a=C_{12}\), \(b=C_{23}\), \(c=C_{34}\). If \[
 a>0,\qquad c>0,\qquad ac>b^2,
\] then the unique optimizer is \[
 X_0=\begin{pmatrix}1&1&0&0\\1&1&0&0\\0&0&1&1\\0&0&1&1\end{pmatrix}.
\] It loses the required edge \(23\). These conditions hold on a
nonempty open set of positive-semidefinite rank-two objectives.

\textbf{Proof.} Put \(q_1=(1,-1,0,0)^T\), \(q_2=(0,0,1,-1)^T\), and \[
 S=[q_1\ q_2]\begin{pmatrix}a&b\\b&c\end{pmatrix}[q_1\ q_2]^T.
\] Then \(S\succeq0\), has rank two, and its entries on the three path
edges are \(-a,-b,-c\). For every feasible \(X\), the quantity
\(\langle C+S,X\rangle\) is constant: all free off-diagonal entries have
zero coefficients in \(C+S\), and the diagonal is fixed. Since
\(\langle S,X\rangle\ge0\), a feasible matrix with \(SX=0\) is optimal.
The matrix \(X_0\) is such a matrix.

Any optimizer must satisfy \(\langle S,X\rangle=0\), hence \(SX=0\).
Thus its range lies in \[
 \ker S=\operatorname{span}\{(1,1,0,0)^T,(0,0,1,1)^T\}.
\] The prescribed unit diagonal forces each diagonal two-by-two block to
be all ones. Every cross-block entry must have the same value; the
constraint \(X_{13}=0\) makes that value zero. This proves uniqueness.

For an explicit rank-two example, take \[
 V=\begin{pmatrix}1&0\\1&1/4\\1/4&1\\1/8&1\end{pmatrix},\qquad C=VV^T.
\] Then \(a=1\), \(b=1/2\), \(c=33/32\), and \(ac-b^2=25/32>0\). The
strict inequalities and full column rank persist under sufficiently
small perturbations of \(V\), so this is an open-set phenomenon within
the rank-two PSD objective family. \(\square\)

In this explicit example \(B=C+S\) has exactly the support of the
complementary path: its nonedge-of-the-path entries are \(B_{13}=3/4\),
\(B_{14}=-3/8\), \(B_{24}=7/8\), and its path-edge entries vanish. Its
rank is four: \(C\) and \(S\) are PSD rank-two matrices whose ranges
together span \(\mathbb R^4\). For example, the two sums of the rows of
\(V\) in the two path-pairs are \((2,1/4)\) and \((3/8,2)\), with
determinant \(125/32\ne0\); this makes the common kernel of \(C,S\)
trivial. Hence the tempting pair has ranks \(2+4=6=n+2\), while the
first matrix is not admissible for the path.

This is a counterexample to the asserted automatic support recovery,
\textbf{not} to SP-10, not to the existence of some better objective,
and not to all semidefinite approaches. An alternative rank-two
objective selection rule would need a separate existence proof for
faithful support on both sides.

\subsection{6. The corrected universal bound and its
limitation}\label{the-corrected-universal-bound-and-its-limitation}

The previously withdrawn statement concerned degeneracies, not ranks.
The false premise \[
 d(G)+d(\overline G)\ge\left\lceil\frac{2n-2}{3}\right\rceil
\] is not used here. The explicit counterexamples in the third archive
remain valid.

One can still recover the numerical rank bound from two different
inputs. Hall's minimum-degree theorem together with minor monotonicity
gives \(\rho(G)\le n-d(G)\). The elementary degeneracy construction
gives \(\rho(\overline G)\le2d(G)+1\). These imply \[
 2\rho(G)+\rho(\overline G)\le2n+1.
\] Swap the graphs and add: \[
 \boxed{\rho(G)+\rho(\overline G)\le\left\lfloor\frac{4n+2}{3}\right\rfloor.}
\] This subsection imports Hall's new theorem {[}1{]}; Sections 2, 3,
and 5 do not. The original proof of the elementary degeneracy
construction is retained in the preceding reports: in dimension
\(2d+1\), keep every set of at most \(d+1\) vectors independent and
follow a degeneracy ordering.

The Strong-Arnold refinement investigated in the interrupted
continuation requires the stronger faithful-SAP version of the
degeneracy bound, attributed there to Mitchell. When that external input
is used, it should be stated at its correct parameter:
\(\nu(\overline G)\ge n-2d(G)-1\), rather than the false degeneracy-sum
inequality. Together with Hall it yields \[
 2\nu(G)+\nu(\overline G)\ge n-1,
\quad
 \nu(G)+\nu(\overline G)\ge\left\lceil\frac{2n-2}{3}\right\rceil.
\] The source audit and exact bibliographic record, when present among
the recovered fourth-continuation files, are preserved rather than
guessed. None of the new class proofs relies on this additional
refinement.

Even the valid universal upper bound has coefficient \(4/3\), not one.
At \(n=100\) it gives 134 rather than 102. Those numbers describe the
gap in the estimate, not a graph violating SP-10.

\subsection{7. Exact supplemental
checks}\label{exact-supplemental-checks}

The supplemental code generates polynomial complement witnesses for all
labeled bipartite patterns with two prescribed parts of size three, for
the explicit half-graph examples, and for a larger one-sided degree-four
fixture. Every support condition is checked using rational arithmetic,
and every recorded rank is recomputed by elimination. A separate unit
test checks the defective path SDP pair: the complement matrix is
accepted, while the primal matrix is required to fail at its missing
middle edge. A test expecting rejection is a soundness test, not a
successful graph-complement certificate.

The proofs, not the finite fixtures, establish the all-order theorems.
The polynomial files certify the complement alone unless an explicit
admissible matrix for the other graph is separately supplied. In
particular this report does not reclassify them as successful rank-pair
certificates.

Current execution results are written to
\nolinkurl{results/supplemental_checks.json}. Recovered logs are stored
separately with the recovered files. Their historical tests are not
silently reported as rerun by the supplemental program.

\subsection{8. What is still needed for a complete
solution}\label{what-is-still-needed-for-a-complete-solution}

For an arbitrary graph one must construct two exactly supported
symmetric matrices of total rank at most \(n+2\), or prove the
equivalent maximum-nullity inequality. No structural reduction to the
degree-four bipartite family has been established. Arbitrary graphs are
not necessarily bipartite; a bipartite graph can also have degrees at
least five in both parts.

The SDP rank count removes neither difficulty: exact support can fail on
an open set, even for a four-vertex path. The successful module theorems
from preceding rounds likewise do not decompose an arbitrary graph into
the required modules. Finally, the valid universal estimate has the
wrong coefficient. These are precise mathematical gaps; no claim of a
complete proof, counterexample, formal proof-assistant verification, or
independent referee approval is made.

\subsection{References and provenance}\label{references-and-provenance}

{[}1{]} H. Tracy Hall, \emph{The Delta Theorem: A dimension bound for
faithful orthogonal graph representations}, arXiv:2601.01211v1, January
2026. \url{https://arxiv.org/html/2601.01211v1} . The minor-monotonicity
background and the newer Delta Theorem are separate dependencies in this
report.

{[}2{]} F. Barioli, S. M. Fallat, H. Gupta, Z. Li, \emph{The weak
version of the graph complement conjecture and partial results for the
delta conjecture}, arXiv:2505.24577v1.
\url{https://arxiv.org/html/2505.24577v1} . Background on the degeneracy
approach; not cited as a proof of the full conjecture.

{[}3{]} C. Erickson, L. Gan, J. Kritschgau, J. C.-H. Lin, S. Spiro,
\emph{Complementary Vanishing Graphs}, arXiv:2207.07294v1.
\url{https://arxiv.org/html/2207.07294v1} . Background on complementary
matrices and the importance of not identifying auxiliary properties with
GCC.

{[}4{]} User-specified SP-10 repository:
\url{https://github.com/ajt60gaibb/OpenProblemsInNLA/tree/main/eigenvalues-and-inverse-problems/SP-10}
. This report does not assert an unverified repository status or modify
the repository.

{[}5{]} Prior conversation artifacts:
\nolinkurl{SP10_partial_results.zip}, \nolinkurl{SP10_round2_results.zip},
and \nolinkurl{SP10_round3_research.zip}, preserved when available. The
third archive corrects the erroneous degeneracy premise in the second.
See \nolinkurl{provenance/recovery.json} for the inventory of the recovered
interrupted continuation.

\end{document}
